\documentclass[12pt]{article}
\usepackage{newtxtext,newtxmath} % Times for text and math
\usepackage{newtxtt} % Modern typewriter for code
\usepackage{listings}
\usepackage{amsmath}
\usepackage{graphicx}
\usepackage{algorithm}
\usepackage{refcount}
\usepackage{hyperref}
\usepackage{xcolor}

\definecolor{codegray}{rgb}{0.5,0.5,0.5}
\definecolor{codegreen}{rgb}{0,0.6,0}
\definecolor{codepurple}{rgb}{0.58,0,0.82}
\definecolor{backcolor}{rgb}{0.95,0.95,0.92}

\lstdefinestyle{mystyle}{
    backgroundcolor=\color{backcolor},   
    commentstyle=\color{codegreen}\itshape,
    keywordstyle=\color{blue}\bfseries,
    numberstyle=\tiny\color{codegray},
    stringstyle=\color{codepurple},
    basicstyle=\ttfamily\footnotesize,
    breaklines=true,
    numbers=left,
    numbersep=5pt,
    showstringspaces=false,
    tabsize=4,
    frame=single,
    captionpos=b
}

\lstset{style=mystyle}

\usepackage{datetime}
\newdateformat{monthdayyeardate}{\monthname[\THEMONTH]~\THEDAY, \THEYEAR}

\begin{document}

\title{\textbf{\LARGE TorchDiff: A PyTorch Library for Denoising Diffusion Probabilistic Models and Beyond}}
\author{Loghman Samani \\ \href{mailto:samaniloqman91@gmail.com}{samaniloqman91@gmail.com} \\ Developer of TorchDiff}
\date{\monthdayyeardate\today}
\maketitle

\begin{abstract}
Diffusion models are a powerful class of generative models that iteratively transform complex data distributions into simple noise distributions and learn to reverse this process to generate new samples. \href{https://github.com/LoqmanSamani/TorchDiff}{TorchDiff} is an open-source Python library built on PyTorch, designed to implement popular diffusion models, including Denoising Diffusion Probabilistic Models (DDPM), Denoising Diffusion Implicit Models (DDIM), Score-Based Generative Models (SDE), and Latent Diffusion Models (LDM). This article introduces the theoretical foundations of diffusion models, focusing on DDPM, and demonstrates how to implement them using TorchDiff. We provide practical examples and invite contributions to the library's development.
\end{abstract}

\section{Introduction to TorchDiff}

\href{https://github.com/LoqmanSamani/TorchDiff}{TorchDiff} is an open-source Python library built on the PyTorch API, designed to facilitate the implementation of state-of-the-art diffusion models. A C++ version is planned for future release to enhance performance. The library is grounded in seminal research, including:

\begin{itemize}
    \item Denoising Diffusion Probabilistic Models (DDPM) by Ho et al. (2020),
    \item Denoising Diffusion Implicit Models (DDIM) by Song et al. (2021),
    \item Score-Based Generative Modeling through Stochastic Differential Equations (SDE) by Song et al. (2021), supporting Variance Exploding (VE), Variance Preserving (VP), and sub-Variance Preserving (sub-VP), and ODE methods.
    \item High-Resolution Image Synthesis with Latent Diffusion Models (LDM) by Rombach et al. (2022).
\end{itemize}

Version 1.0.0 of TorchDiff supports both conditional and unconditional training and sampling for these models. It includes a transformer-based text encoder, enabling conditioning on text prompts or class labels, making it suitable for applications like text-to-image generation. Readers are encouraged to visit the \href{https://github.com/LoqmanSamani/TorchDiff}{TorchDiff GitHub repository}, explore the documentation at \href{https://torchdiff.readthedocs.io}{torchdiff.readthedocs.io}, and contribute to its development.

\section{Introduction to Diffusion Models}

Diffusion models are a class of deep generative models that generate data by learning to reverse a process that gradually adds noise to a complex data distribution. Introduced by Sohl-Dickstein et al. (2015), the concept draws inspiration from non-equilibrium statistical physics, specifically the Jarzynski equality, which relates non-equilibrium processes to equilibrium free energy differences \cite{jarzynski1997}. The core idea is to iteratively corrupt a data distribution (e.g., images) through a forward process, transforming it into a simple noise distribution, and then learn to reverse this process to generate new samples (Figure \ref{fig:diffusion}).

\begin{figure}[htbp]
    \centering
    \includegraphics[width=1\textwidth]{diff1.png}
    \caption{Forward and reverse processes of a diffusion model trained on 2D Swiss roll data. The top row shows the forward process $q(x^{(0...T)})$, where the data distribution (left) is gradually corrupted into an isotropic Gaussian (right). The bottom row shows the reverse process $p(x^{(0...T)})$, reconstructing the data distribution. Source: Sohl-Dickstein et al. (2015).}
    \label{fig:diffusion}
\end{figure}

\subsection{Historical Context}

Historically, probabilistic generative models have faced a trade-off between tractability and flexibility. Tractable models, such as Gaussian distributions, are easy to analyze and fit but lack the expressiveness to capture complex patterns in high-dimensional data (e.g., images). Flexible models, such as those requiring Monte Carlo sampling, can model intricate distributions but are computationally expensive to train and evaluate \cite{neal2001}. Sohl-Dickstein et al. (2015) proposed diffusion probabilistic models to address this trade-off, offering both analytical tractability and the ability to model complex data distributions.

In their seminal work on DDPM, Ho et al. (2020) refined this approach by using a Markov chain to gradually transform a complex data distribution into a standard Gaussian distribution (Figure \ref{fig:ddpm}). By ensuring small noise increments at each step, the forward process becomes analytically tractable, enabling efficient training and high-quality sample generation.

\begin{figure}[htbp]
    \centering
    \includegraphics[width=1\textwidth]{diff2.png}
    \caption{Forward and reverse processes of DDPM. The forward process (denoted by $q(x_t|x_{t-1})$, dashed arrows) gradually adds noise to the original image $x_0$ until it becomes pure noise $x_T$. The reverse process (denoted by $p_\theta(x_{t-1} | x_t)$, solid arrows) is learned to gradually denoise the sample from $x_T$ back to $x_0$, reconstructing the original image. Source: Ho et al. (2020).}
    \label{fig:ddpm}
\end{figure}

\section{Forward Process of Diffusion Models}
\label{sec:forward}

The forward process in a diffusion model systematically transforms a complex data distribution, denoted \( q(x^{(0)}) \), into a simple, tractable distribution, typically an isotropic Gaussian \( \pi(x^{(T)}) \). This is achieved through a sequence of Markov diffusion kernels, each adding a small amount of Gaussian noise to the data. The process is inspired by non-equilibrium statistical physics, where structured information is gradually eroded, much like adding static to a radio signal until only noise remains \cite{sohl2015}.

Formally, the forward process is defined as a Markov chain over \( T \) steps, with the joint probability:

\begin{align}
q(x^{(0:T)}) = q(x^{(0)}) \prod_{t=1}^{T} q(x^{(t)} | x^{(t-1)})
\end{align}

Each transition \( q(x^{(t)} | x^{(t-1)}) \), known as a Markov diffusion kernel, is a Gaussian distribution that slightly attenuates the signal and adds noise:

\begin{align}
q(x^{(t)} | x^{(t-1)}) = \mathcal{N}(x^{(t)}; \sqrt{1 - \beta_t} x^{(t-1)}, \beta_t I)
\end{align}

Here, \( x^{(t-1)} \) is the data at step \( t-1 \), \( \sqrt{1 - \beta_t} \) scales the signal to preserve its magnitude, and \( \beta_t I \) introduces Gaussian noise with variance \( \beta_t \). The parameter \( \beta_t \), known as the noise schedule, controls the amount of noise added at each step. Over \( T \) steps, this process transforms the data distribution into pure noise, approximating \( \pi(x^{(T)}) \approx \mathcal{N}(0, I) \).

A key property of the forward process is its analytical tractability. For small \( \beta_t \), the transitions are nearly deterministic, allowing the full chain to be evaluated efficiently. Additionally, the forward process can be expressed directly from the initial data \( x^{(0)} \) to any step \( t \) using a reparameterization:

\begin{align}
x^{(t)} = \sqrt{\bar{\alpha}_t} x^{(0)} + \sqrt{1 - \bar{\alpha}_t} \epsilon, \quad \epsilon \sim \mathcal{N}(0, I)
\end{align}

where \( \bar{\alpha}_t = \prod_{s=1}^{t} (1 - \beta_s) \). This formulation, introduced by Ho et al. (2020), enables efficient sampling during training by directly computing \( x^{(t)} \) without iterating through all intermediate steps \cite{ho2020}.

\subsection{Noise Schedule (\(\beta_t\))}
\label{subsec:noise_schedule}

The noise schedule \( \beta_t \) is a critical hyperparameter that determines the rate at which noise is added. In the original diffusion model by Sohl-Dickstein et al. (2015), \( \beta_t \) was treated as a trainable parameter, optimized during training \cite{sohl2015}. However, Ho et al. (2020) found that a fixed noise schedule, such as a linear or cosine schedule, simplifies training and improves sample quality \cite{ho2020}. A fixed schedule offers several advantages:

\begin{itemize}
    \item \textbf{Stability}: A predefined schedule reduces the complexity of optimizing \( \beta_t \), leading to more stable training dynamics.
    \item \textbf{Simplicity}: It eliminates the need for additional hyperparameters or learning algorithms for \( \beta_t \).
    \item \textbf{Quality}: Carefully designed schedules, such as linear or cosine, ensure gradual noise addition, preserving structural information in early steps and achieving pure noise by step \( T \).
\end{itemize}

In a linear schedule, \( \beta_t \) increases linearly from a small value (e.g., \( 10^{-4} \)) to a larger value (e.g., \( 0.02 \)) over \( T \) steps. This gradual increase ensures that early steps preserve the data’s structure, allowing the model to learn meaningful features, while later steps produce near-pure noise. Table \ref{tab:beta_schedule} illustrates a linear noise schedule for a DDPM model.

\begin{table}[htbp]
    \centering
    \caption{Example linear noise schedule for DDPM with \( T = 1000 \), \( \beta_1 = 10^{-4} \), and \( \beta_T = 0.02 \).}
    \label{tab:beta_schedule}
    \begin{tabular}{c|c|c|c}
        \hline
        Step \( t \) & \( \beta_t \) & \( \alpha_t = 1 - \beta_t \) & \( \bar{\alpha}_t = \prod_{s=1}^{t} \alpha_s \) \\
        \hline
        1 & \( 10^{-4} \) & 0.9999 & 0.9999 \\
        500 & \( 0.01 \) & 0.99 & 0.0821 \\
        1000 & 0.02 & 0.98 & 0.0002 \\
        \hline
    \end{tabular}
\end{table}

\subsection{Initializing the Forward Process in TorchDiff}
\label{subsec:torchdiff_forward}

The \href{https://torchdiff.readthedocs.io}{TorchDiff} library, built on PyTorch, provides a modular interface for implementing diffusion models, including DDPM. To initialize the forward process for DDPM, users first define a set of hyperparameters using the \texttt{HyperParamsDDPM} class, which specifies the number of diffusion steps and the noise schedule. The following code snippet demonstrates this process:

\begin{lstlisting}[language=Python, caption={Initializing DDPM hyperparameters in TorchDiff.}, label={lst:hyperparams}]
import torch
from torchdiff.ddpm import HyperParamsDDPM

# Define DDPM hyperparameters
hp_ddpm = HyperParamsDDPM(
    num_steps=1000,          # Number of diffusion steps
    beta_start=1e-4,         # Initial noise scale
    beta_end=0.02,           # Final noise scale
    beta_method="linear"     # Linear noise schedule
)
\end{lstlisting}

In Listing \ref{lst:hyperparams}, \texttt{num\_steps} determines the length of the forward process, during which the data is corrupted into Gaussian noise. The parameters \texttt{beta\_start} and \texttt{beta\_end} define the range of the noise schedule, starting with a small \( \beta_1 = 10^{-4} \) to preserve data structure and increasing to \( \beta_T = 0.02 \) to ensure pure noise at step \( T \). The \texttt{beta\_method="linear"} argument specifies a linear interpolation for \( \beta_t \), though TorchDiff also supports alternatives like cosine or sigmoid schedules for advanced use cases.

Next, the \texttt{ForwardDDPM} class is initialized using the defined hyperparameters to set up the forward diffusion process:

\begin{lstlisting}[language=Python, caption={Initializing the forward process in TorchDiff.}, label={lst:forward_ddpm}]
from torchdiff.ddpm import ForwardDDPM

# Initialize forward diffusion process
fwd_ddpm = ForwardDDPM(hp_ddpm)
\end{lstlisting}

The \texttt{ForwardDDPM} class encapsulates the logic for applying the Gaussian diffusion kernels, computing intermediate states \( x^{(t)} \), and managing the noise schedule. This modular design allows users to easily integrate the forward process into training and sampling pipelines, as demonstrated in subsequent sections.




\section{Reverse Process of Diffusion Models}
\label{sec:reverse}

The reverse process in a diffusion model is the generative counterpart to the forward process, aiming to reconstruct the original data distribution \( q(x^{(0)}) \) from pure Gaussian noise \( \pi(x^{(T)}) \). Starting with a sample \( x^{(T)} \sim \mathcal{N}(0, I) \), the model iteratively applies learned reverse transitions \( p_\theta(x^{(t-1)} | x^{(t)}) \) to denoise the data, gradually recovering the structure of the original distribution. This process is akin to restoring a blurred image by iteratively sharpening it, guided by a learned model.

Formally, the reverse process is defined as a Markov chain, with the joint probability:

\begin{align}
p_\theta(x^{(0:T)}) = p(x^{(T)}) \prod_{t=1}^{T} p_\theta(x^{(t-1)} | x^{(t)})
\end{align}

where \( p(x^{(T)}) = \mathcal{N}(0, I) \) is the initial noise distribution, and each reverse transition is a Gaussian diffusion kernel parameterized by a neural network:

\begin{align}
p_\theta(x^{(t-1)} | x^{(t)}) = \mathcal{N}(x^{(t-1)}; \mu_\theta(x^{(t)}, t), \Sigma_\theta(x^{(t)}, t))
\end{align}

Here, \( \mu_\theta(x^{(t)}, t) \) is the mean of the reverse transition, and \( \Sigma_\theta(x^{(t)}, t) \) is the covariance, both predicted by a neural network parameterized by \( \theta \). In the original diffusion model by Sohl-Dickstein et al. (2015), both the mean and covariance were learned \cite{sohl2015}. However, in the DDPM framework by Ho et al. (2020), the covariance \( \Sigma_\theta \) is fixed to a predefined schedule (e.g., \( \Sigma_\theta = \beta_t I \)), and the neural network focuses on predicting the mean \( \mu_\theta \) \cite{ho2020}.

A key insight in DDPM is that instead of directly predicting the mean, the neural network, referred to as the \emph{noise predictor}, is trained to estimate the noise component \( \epsilon \) added during the forward process. This is mathematically equivalent because the mean \( \mu_\theta \) can be derived from the predicted noise using the forward process parameters. Specifically, for a sample \( x^{(t)} = \sqrt{\bar{\alpha}_t} x^{(0)} + \sqrt{1 - \bar{\alpha}_t} \epsilon \), the noise predictor \( \epsilon_\theta(x^{(t)}, t) \) estimates \( \epsilon \), and the mean is computed as:

\begin{align}
\mu_\theta(x^{(t)}, t) = \frac{1}{\sqrt{\alpha_t}} \left( x^{(t)} - \frac{\beta_t}{\sqrt{1 - \bar{\alpha}_t}} \epsilon_\theta(x^{(t)}, t) \right)
\end{align}

where \( \alpha_t = 1 - \beta_t \) and \( \bar{\alpha}_t = \prod_{s=1}^{t} \alpha_s \). This formulation simplifies training and leverages the forward process’s tractability, as discussed in Section \ref{sec:forward}.


\subsection{Noise Predictor Architecture}
\label{subsec:noise_predictor}

In DDPM, the noise predictor is typically a U-Net-like neural network, which incorporates time embeddings and self-attention mechanisms to model the denoising process effectively \cite{ho2020}. The U-Net architecture consists of downsampling, bottleneck, and upsampling blocks, allowing it to capture both local and global features in the data (e.g., image textures and structures). Time embeddings, which encode the diffusion step \( t \), are injected into the network to make the predictions time-conditional. Multi-head self-attention layers enhance the model’s ability to focus on relevant parts of the input, improving the quality of generated samples.

The noise predictor takes as input a noisy sample \( x^{(t)} \) and the time step \( t \), outputting the estimated noise \( \epsilon_\theta(x^{(t)}, t) \). During training, the model learns to minimize the difference between the predicted noise and the actual noise added in the forward process, as detailed in Section \ref{subsec:training}.

\subsection{Initializing the Reverse Process in TorchDiff}
\label{subsec:torchdiff_reverse}

The \href{https://torchdiff.readthedocs.io}{TorchDiff} library provides a straightforward interface for initializing the reverse process of DDPM. The \texttt{ReverseDDPM} class is used to set up the reverse diffusion process, leveraging the same hyperparameters defined for the forward process (Section \ref{subsec:torchdiff_forward}). The following code snippet demonstrates this initialization:

\begin{lstlisting}[language=Python, caption={Initializing the reverse process in TorchDiff.}, label={lst:reverse_ddpm}]
from torchdiff.ddpm import ReverseDDPM

# Initialize reverse diffusion process
rev_ddpm = ReverseDDPM(hp_ddpm)
\end{lstlisting}

In Listing \ref{lst:reverse_ddpm}, the \texttt{ReverseDDPM} class uses the \texttt{hp\_ddpm} object, which contains the noise schedule (\( \beta_t \), \( \alpha_t \), \( \bar{\alpha}_t \)) and other parameters computed during forward process initialization. These parameters ensure consistency between the forward and reverse processes, enabling accurate denoising.

To implement the noise predictor, TorchDiff provides the \texttt{NoisePredictor} class, which mirrors the U-Net architecture described in Ho et al. (2020). The following code snippet initializes a noise predictor for RGB images:

\begin{lstlisting}[language=Python, caption={Initializing the noise predictor in TorchDiff.}, label={lst:noise_predictor}]
from torchdiff.utils import NoisePredictor

# Initialize noise predictor model
noise_pred = NoisePredictor(
    in_channels=3,              # Input channels (3 for RGB images)
    down_channels=[32, 64, 128, 256],  # Channels for downsampling blocks
    mid_channels=[256, 256, 256, 256], # Channels for bottleneck blocks
    up_channels=[256, 128, 64, 32],    # Channels for upsampling blocks
    time_embed_dim=256,         # Dimension of time embeddings
    y_embed_dim=256,            # Dimension of conditional embeddings (e.g., text)
    num_down_blocks=2,          # Number of downsampling blocks
    num_mid_blocks=2,           # Number of bottleneck blocks
    num_up_blocks=2             # Number of upsampling blocks
)
\end{lstlisting}

In Listing \ref{lst:noise_predictor}, the \texttt{NoisePredictor} is configured with a U-Net architecture tailored for RGB images (\texttt{in\_channels=3}). The \texttt{down\_channels}, \texttt{mid\_channels}, and \texttt{up\_channels} define the channel progression through the U-Net’s encoder, bottleneck, and decoder, respectively. The \texttt{time\_embed\_dim} parameter specifies the size of the time embeddings, enabling the model to condition predictions on the diffusion step \( t \). The \texttt{y\_embed\_dim} parameter supports conditional generation (e.g., text prompts), which is discussed further in Section \ref{sec:training}. The number of blocks (\texttt{num\_down\_blocks}, \texttt{num\_mid\_blocks}, \texttt{num\_up\_blocks}) controls the depth of the network, balancing computational complexity and model capacity.



\section{Training Objective}
\label{sec:training}

The training objective of a diffusion model is to optimize the model parameters \( \theta \) to maximize the likelihood of the observed data under the generative model. This involves computing the probability \( p_\theta(x^{(0)}) \) that the model assigns to real data points \( x^{(0)} \), measuring how well the model captures the true data distribution. This section outlines the general approach to maximizing log-likelihood, the KL divergence-based loss formulation from Sohl-Dickstein et al. (2015), and the simplified loss introduced by Ho et al. (2020) for DDPM \cite{sohl2015,ho2020}.


\subsection{Maximizing Log-Likelihood}
\label{subsec:log_likelihood}

A key question in training diffusion models is how to evaluate the model’s fit to the data. The goal is to maximize the expected log-likelihood over the true data distribution \( q(x^{(0)}) \):

\begin{align}
\mathcal{L} = \mathbb{E}_{q(x^{(0)})} \left[ \log p_\theta(x^{(0)}) \right]
\end{align}

The generative model defines the probability \( p_\theta(x^{(0)}) \) through the joint distribution over the full trajectory:

\begin{align}
p_\theta(x^{(0)}) = \int dx^{(1:T)} \, p_\theta(x^{(0:T)}) = \int dx^{(1:T)} \, p(x^{(T)}) \prod_{t=1}^{T} p_\theta(x^{(t-1)} | x^{(t)})
\end{align}

However, computing this integral over all possible noise trajectories is intractable due to the high-dimensional nature of \( x^{(1:T)} \). To address this, diffusion models leverage the forward process \( q(x^{(1:T)} | x^{(0)}) \), which is analytically tractable, as discussed in Section \ref{sec:forward}. Using importance sampling, the likelihood can be rewritten as:

\begin{align}
p_\theta(x^{(0)}) = \int dx^{(1:T)} \, q(x^{(1:T)} | x^{(0)}) \, \frac{p_\theta(x^{(0:T)})}{q(x^{(1:T)} | x^{(0)})}
\end{align}

This allows estimation of \( p_\theta(x^{(0)}) \) by sampling trajectories from the forward process and computing the ratio of reverse and forward transition probabilities:

\begin{align}
p_\theta(x^{(0)}) \approx p(x^{(T)}) \prod_{t=1}^{T} \frac{p_\theta(x^{(t-1)} | x^{(t)})}{q(x^{(t)} | x^{(t-1)})}
\end{align}

When the noise steps \( \beta_t \) are small, the forward and reverse processes are nearly symmetric, and this approximation becomes accurate with a single trajectory sample, resembling a quasi-static process inspired by the Jarzynski equality \cite{jarzynski1997}.

To make training tractable, the log-likelihood is bounded using Jensen’s inequality, yielding the evidence lower bound (ELBO):

\begin{align}
\log p_\theta(x^{(0)}) \geq \mathbb{E}_{q(x^{(1:T)} | x^{(0)})} \left[ \log \frac{p_\theta(x^{(0:T)})}{q(x^{(1:T)} | x^{(0)})} \right]
\end{align}

Taking the expectation over \( q(x^{(0)}) \), the training objective becomes maximizing the ELBO:

\begin{align}
\mathcal{L}_{\text{ELBO}} = \mathbb{E}_{q(x^{(0:T)})} \left[ \log \frac{p_\theta(x^{(0:T)})}{q(x^{(1:T)} | x^{(0)})} \right]
\end{align}

This ELBO measures how closely the learned reverse process \( p_\theta \) matches the true posterior \( q \), providing a practical objective for optimization.

\subsection{KL Formulation of the Loss (2015)}
\label{subsec:kl_2015}

Sohl-Dickstein et al. (2015) decomposed the ELBO into a sum of Kullback-Leibler (KL) divergences and entropy terms, offering a theoretically grounded loss function \cite{sohl2015}. The KL divergence measures the difference between two probability distributions, here comparing the true posterior \( q(x^{(t-1)} | x^{(t)}, x^{(0)}) \) to the learned reverse transition \( p_\theta(x^{(t-1)} | x^{(t)}) \). The ELBO can be expressed as:

\begin{align}
\mathcal{L}_{\text{ELBO}} = & - \sum_{t=2}^{T} \mathbb{E}_{q(x^{(0)}, x^{(t)})} \left[ 
    D_{\text{KL}} \left( q(x^{(t-1)} | x^{(t)}, x^{(0)}) \parallel p_\theta(x^{(t-1)} | x^{(t)}) \right) 
\right] \nonumber \\
& + H_q(x^{(T)} | x^{(0)}) - H_q(x^{(1)} | x^{(0)}) - H_p(x^{(T)})
\end{align}


where:
\begin{itemize}
    \item \( D_{\text{KL}}(q \parallel p) \): The KL divergence between the true and learned reverse transitions at each step \( t \).
    \item \( H_q(x^{(t)} | x^{(0)}) \): Conditional entropy terms capturing uncertainty in the forward process.
    \item \( H_p(x^{(T)}) \): Entropy of the initial noise distribution, typically a constant for \( p(x^{(T)}) = \mathcal{N}(0, I) \).
\end{itemize}

Training involves minimizing the sum of KL divergences, effectively aligning the reverse process with the true data-generating process. This formulation is computationally intensive due to the need to estimate multiple KL terms, but it provides a robust framework for learning flexible diffusion models.

\subsection{Simplified DDPM Loss (2020)}
\label{subsec:ddpm_loss}

Ho et al. (2020) introduced a simplified loss function for DDPM that significantly reduces computational complexity while maintaining high sample quality \cite{ho2020}. They proposed three key simplifications:
\begin{itemize}
    \item \textbf{Fixed Noise Schedule}: The forward process noise schedule \( \beta_t \) is predefined (e.g., linear or cosine), as discussed in Section \ref{subsec:noise_schedule}.
    \item \textbf{Fixed Reverse Variance}: The reverse process covariance \( \Sigma_\theta(x^{(t)}, t) \) is set to a constant (e.g., \( \beta_t I \)) rather than learned.
    \item \textbf{Noise Prediction}: The neural network predicts the noise \( \epsilon \) added in the forward process instead of the reverse mean \( \mu_\theta \).
\end{itemize}

Recall from Section \ref{sec:forward} that the forward process generates a noisy sample:

\begin{align}
x^{(t)} = \sqrt{\bar{\alpha}_t} x^{(0)} + \sqrt{1 - \bar{\alpha}_t} \epsilon, \quad \epsilon \sim \mathcal{N}(0, I)
\end{align}

where \( \bar{\alpha}_t = \prod_{s=1}^{t} (1 - \beta_s) \). The noise predictor \( \epsilon_\theta(x^{(t)}, t) \) is trained to estimate \( \epsilon \) given \( x^{(t)} \) and \( t \). The simplified loss is the mean squared error (MSE) between the true and predicted noise:

\begin{align}
L_{\text{simple}}(\theta) = \mathbb{E}_{t, x^{(0)}, \epsilon} \left[ \left\| \epsilon - \epsilon_\theta \left( \sqrt{\bar{\alpha}_t} x^{(0)} + \sqrt{1 - \bar{\alpha}_t} \epsilon, t \right) \right\|_2^2 \right]
\end{align}

This loss approximates the ELBO because predicting the noise \( \epsilon \) is equivalent to predicting the reverse mean \( \mu_\theta \), as shown in Section \ref{sec:reverse}. The simplification reduces the training objective to a straightforward regression task, making it computationally efficient and empirically effective. Ho et al. (2020) demonstrated that this loss produces high-quality samples, comparable to or better than the full KL-based loss, due to its focus on denoising performance.

Table \ref{tab:loss_comparison} summarizes the differences between the 2015 and 2020 loss formulations.

\begin{table*}[htbp]
    \centering
    \caption{Comparison of loss formulations for diffusion models.}
    \label{tab:loss_comparison}
    \small % or \scriptsize to reduce further
    \begin{tabular}{p{3.5cm}|p{6cm}|p{6cm}}
        \hline
        \textbf{Aspect} & \textbf{Sohl-Dickstein et al. (2015)} & \textbf{Ho et al. (2020)} \\
        \hline
        \textbf{Objective} & Sum of KL divergences (ELBO) & Simplified MSE on noise prediction \\
        \textbf{Noise Schedule} & Trainable \( \beta_t \) & Fixed (e.g., linear, cosine) \\
        \textbf{Reverse Variance} & Learned \( \Sigma_\theta \) & Fixed (e.g., \( \beta_t I \)) \\
        \textbf{Prediction Target} & Mean and variance of reverse transition & Noise component \( \epsilon \) \\
        \textbf{Computational Complexity} & High (multiple KL terms) & Low (single MSE term) \\
        \hline
    \end{tabular}
\end{table*}



\subsection{Training with TorchDiff}
\label{subsec:torchdiff_training}

With the theoretical foundations of DDPM established, this section demonstrates how to train a conditional DDPM model using the \href{https://torchdiff.readthedocs.io}{TorchDiff} library. Training involves preparing a dataset, configuring a conditional text encoder, defining evaluation metrics, and setting up the training pipeline. The following steps outline the process, focusing on a practical example with ImageNet data.

First, the dataset is prepared using PyTorch’s \texttt{torchvision} utilities to load and preprocess images. The images are normalized to the range \([-1, 1]\), which is standard for diffusion models to ensure stable training:\\[5px]

\begin{lstlisting}[language=Python, caption={Preparing the ImageNet dataset for training.}, label={lst:dataset}]
from torchvision import datasets, transforms
from torch.utils.data import DataLoader

# Define image preprocessing transformations
transform = transforms.Compose([
    transforms.Resize(224),           # Resize to 224x224 pixels
    transforms.CenterCrop(224),       # Crop to center
    transforms.ToTensor(),            # Convert to tensor
    transforms.Normalize((0.5, 0.5, 0.5), (0.5, 0.5, 0.5))  # Normalize to [-1, 1]
])

# Load ImageNet training and validation datasets
train_dataset = datasets.ImageFolder(root='./imagenet/train', transform=transform)
val_dataset = datasets.ImageFolder(root='./imagenet/val', transform=transform)

# Create data loaders
train_loader = DataLoader(
    train_dataset, batch_size=64, shuffle=True, num_workers=8, pin_memory=True
)
val_loader = DataLoader(
    val_dataset, batch_size=64, shuffle=False, num_workers=8, pin_memory=True
)
\end{lstlisting}

\vspace{3em}

In Listing \ref{lst:dataset}, the \texttt{transform} pipeline resizes and normalizes images, while the \texttt{DataLoader} objects facilitate efficient batch processing. The \texttt{pin\_memory=True} option optimizes data transfer to GPUs, enhancing training performance.

Next, a conditional text encoder is initialized to enable text-prompt-based generation, such as generating images from descriptions like “a cat” or “a dog.” TorchDiff’s \texttt{TextEncoder} class leverages a pre-trained BERT model to encode text prompts into embeddings:
\vspace{3em}

\begin{lstlisting}[language=Python, caption={Initializing the text encoder for conditional generation.}, label={lst:text_encoder}]
from torchdiff.utils import TextEncoder

# Initialize text encoder
text_enc = TextEncoder(
    use_pretrained_model=True,      # Use pre-trained BERT
    model_name="bert-base-uncased", # BERT model variant
    output_dimension=256            # Embedding dimension
)
\end{lstlisting}

In Listing \ref{lst:text_encoder}, the \texttt{TextEncoder} produces 256-dimensional embeddings, which are used to condition the noise predictor during training and sampling. This enables the model to generate images that align with specific text prompts, enhancing its flexibility for applications like text-to-image synthesis.

To evaluate the model’s performance during training, TorchDiff provides the \texttt{Metrics} class, which supports several image quality metrics, including Mean Squared Error (MSE), Peak Signal-to-Noise Ratio (PSNR), Structural Similarity Index (SSIM), Fréchet Inception Distance (FID), and Learned Perceptual Image Patch Similarity (LPIPS). These metrics compare generated images to ground-truth images, providing quantitative insights into model performance:

\begin{lstlisting}[language=Python, caption={Configuring evaluation metrics.}, label={lst:metrics}]
from torchdiff.utils import Metrics

# Initialize metrics for evaluation
metrics = Metrics(
    device="cuda",      # Use GPU for computations
    fid=True,          # Compute FID score
    ssim=True,         # Compute MSE, PSNR, and SSIM
    lpips_=True        # Compute LPIPS using VGG backbone
)
\end{lstlisting}

In Listing \ref{lst:metrics}, \texttt{fid=True} enables FID computation, which measures the similarity between generated and real image distributions using a pre-trained Inception V3 model. Similarly, \texttt{ssim=True} and \texttt{lpips\_=True} activate structural and perceptual similarity metrics, respectively.

Finally, the \texttt{TrainDDPM} class orchestrates the training process, combining the noise predictor, text encoder, hyperparameters, and evaluation metrics. The training objective is the simplified MSE loss discussed in Section \ref{subsec:ddpm_loss}, optimized using the Adam optimizer:

\begin{lstlisting}[language=Python, caption={Setting up and starting DDPM training.}, label={lst:train_ddpm}]
import torch.nn as nn
from torchdiff.ddpm import TrainDDPM

# Configure optimizer
optim = torch.optim.Adam(
    list(noise_pred.parameters()) + list(text_enc.parameters()),  # Optimize both models
    lr=1e-4                                              # Learning rate
)

# Define MSE loss function
loss_fn = nn.MSELoss()

# Initialize training pipeline
trainer = TrainDDPM(
    noise_predictor=noise_pred,     # Noise predictor model
    hyper_params=hp_ddpm,           # Hyperparameters
    conditional_model=text_enc,     # Text encoder for conditioning
    metrics_=metrics,               # Evaluation metrics
    optimizer=optim,                # Adam optimizer
    objective=loss_fn,              # MSE loss function
    data_loader=train_loader,       # Training data loader
    val_loader=val_loader,          # Validation data loader
    max_epoch=100,                  # Number of training epochs
    device="cuda",                  # Use GPU
    store_path="ddpm_model.pth",    # Path to save model checkpoints
    val_frequency=5                 # Evaluate every 5 epochs
)

# Start training
trainer()
\end{lstlisting}

In Listing \ref{lst:train_ddpm}, the \texttt{TrainDDPM} class integrates all components, training the noise predictor and text encoder to minimize the MSE between predicted and actual noise. The \texttt{val\_frequency=5} parameter triggers validation every five epochs, using the metrics defined earlier to monitor performance. The trained model is saved to \texttt{ddpm\_model.pth} for later use in sampling.

\subsection{Sampling with TorchDiff}
\label{subsec:torchdiff_sampling}

Once trained, the DDPM model can generate new samples, either unconditionally or conditioned on text prompts, using the reverse process described in Section \ref{sec:reverse}. TorchDiff’s \texttt{SampleDDPM} class simplifies this process, allowing users to generate images by specifying the model, text prompts, and output settings. The following code snippet demonstrates generating three conditional images:

\begin{lstlisting}[language=Python, caption={Generating conditional images with TorchDiff.}, label={lst:sample_ddpm}]
from torchdiff.ddpm import SampleDDPM

# Initialize sampler
sampler = SampleDDPM(
    reverse_diffusion=rev_ddpm,     # Reverse process
    noise_predictor=noise_pred,     # Noise predictor model
    image_shape=(224, 224),         # Output image dimensions
    conditional_model=text_enc,     # Text encoder for conditioning
    tokenizer="bert-base-uncased",  # Tokenizer for text prompts
    batch_size=3,                   # Generate three images
    in_channels=3,                  # RGB images
    device="cuda",                  # Use GPU
    output_range=(-1, 1)            # Output range matching training data
)

# Generate images from text prompts
images = sampler(
    conditions=["a cat", "a dog", "a box"],  # Text prompts
    save_images=True,                        # Save generated images
    save_path="ddpm_generated"               # Output directory
)
\end{lstlisting}

In Listing \ref{lst:sample_ddpm}, the \texttt{SampleDDPM} class uses the trained noise predictor and text encoder to generate 224x224 RGB images conditioned on the prompts “a cat,” “a dog,” and “a box.” The \texttt{output\_range=(-1, 1)} ensures the generated images match the training data’s normalization. The \texttt{save\_images=True} option saves the images to the \texttt{ddpm\_generated} directory, facilitating visual inspection.

Table \ref{tab:torchdiff_components} summarizes the key TorchDiff classes used in training and sampling, highlighting their roles and configurations.

\begin{table*}[htbp]
    \centering
    \caption{Key TorchDiff classes for DDPM training and sampling.}
    \label{tab:torchdiff_components}
    \scriptsize
    \begin{tabular}{p{3cm}|p{4.5cm}|p{8.5cm}}
        \hline
        \textbf{Class} & \textbf{Role} & \textbf{Key Parameters} \\
        \hline
        \texttt{HyperParamsDDPM} & Defines diffusion hyperparameters & \texttt{num\_steps}, \texttt{beta\_start}, \texttt{beta\_end}, \texttt{beta\_method} \\
        \texttt{ForwardDDPM} & Implements forward process & \texttt{hyper\_params} \\
        \texttt{ReverseDDPM} & Implements reverse process & \texttt{hyper\_params} \\
        \texttt{NoisePredictor} & Predicts noise for denoising & \texttt{in\_channels}, \texttt{down\_channels}, \texttt{time\_embed\_dim} \\
        \texttt{TextEncoder} & Encodes text prompts & \texttt{model\_name}, \texttt{output\_dimension} \\
        \texttt{Metrics} & Evaluates generated images & \texttt{fid}, \texttt{ssim}, \texttt{lpips\_} \\
        \texttt{TrainDDPM} & Orchestrates training & \texttt{noise\_predictor}, \texttt{conditional\_model}, \texttt{max\_epoch} \\
        \texttt{SampleDDPM} & Generates samples & \texttt{reverse\_diffusion}, \texttt{image\_shape}, \texttt{conditions} \\
        \hline
    \end{tabular}
\end{table*}


\section{Conclusion}
\label{sec:conclusion}

Diffusion models, exemplified by Denoising Diffusion Probabilistic Models (DDPM), offer a powerful framework for generative modeling, balancing tractability and flexibility to produce high-quality samples. This article has provided a comprehensive introduction to diffusion models, covering their theoretical foundations—forward and reverse processes, training objectives—and practical implementation using the \href{https://github.com/LoqmanSamani/TorchDiff}{TorchDiff} library. TorchDiff simplifies the development of DDPM and other diffusion models (e.g., DDIM, SDE, LDM) by providing modular, PyTorch-based classes for configuring, training, and sampling.

Key contributions of TorchDiff include its support for both conditional and unconditional generation, a flexible text encoder for text-to-image tasks, and a suite of evaluation metrics to monitor performance. The library’s design, grounded in seminal research \cite{sohl2015,ho2020,song2021,rombach2022}, makes it accessible to researchers and practitioners alike. Future developments, including a planned C++ implementation, promise to enhance performance and broaden its applicability.

Readers are encouraged to explore TorchDiff at \href{https://github.com/LoqmanSamani/TorchDiff}{github.com/LoqmanSamani/TorchDiff}, contribute to its development (e.g., by adding new models or optimizing code), and leverage its capabilities for generative modeling tasks. By bridging theory and practice, TorchDiff empowers the community to advance the field of diffusion models.


\newpage



\begin{thebibliography}{9}
\bibitem{jarzynski1997}
Jarzynski, C. (1997). Nonequilibrium equality for free energy differences. \emph{Physical Review Letters}, 78(14), 2690.

\bibitem{neal2001}
Neal, R. M. (2001). Annealed importance sampling. \emph{Statistics and Computing}, 11(2), 125--139.

\bibitem{sohl2015}
Sohl-Dickstein, J., Weiss, E. A., Maheswaranathan, N., \& Ganguli, S. (2015). Deep unsupervised learning using nonequilibrium thermodynamics. \emph{International Conference on Machine Learning (ICML)}, 2256--2265.

\bibitem{ho2020}
Ho, J., Jain, A., \& Abbeel, P. (2020). Denoising diffusion probabilistic models. \emph{Advances in Neural Information Processing Systems (NeurIPS)}, 33, 6840--6851.

\bibitem{song2021}
Song, J., Meng, C., \& Ermon, S. (2021). Denoising diffusion implicit models. \emph{International Conference on Learning Representations (ICLR)}.

\bibitem{song2021sde}
Song, Y., Sohl-Dickstein, J., Kingma, D. P., Kumar, A., Ermon, S., \& Poole, B. (2021). Score-based generative modeling through stochastic differential equations. \emph{International Conference on Learning Representations (ICLR)}.

\bibitem{rombach2022}
Rombach, R., Blattmann, A., Lorenz, D., Esser, P., \& Ommer, B. (2022). High-resolution image synthesis with latent diffusion models. \emph{Proceedings of the IEEE/CVF Conference on Computer Vision and Pattern Recognition (CVPR)}, 10684--10695.
\end{thebibliography}

\end{document}
